% ============================================================================
% DRAFT: t3_separation_propositions.tex   (W2 deliverable, T3 landing draft)
% Date: 2026-08-04 / package: 2026-08-05
% Source: paper/current-usenix/t3_principle_separation_propositions_2026-08-04.md
% Status: INDEPENDENT DRAFT. Do NOT \input into main.tex until the PM
%         approves placement and the claim-map diff review has run.
%         This file does not modify any existing section.
%
% ----------------------------------------------------------------------------
% 落位建议 (placement suggestion)
%   Part 1 — the two propositions:
%   File    : sections/security_analysis.tex
%   Position: inside \subsection{The Cost of an Insufficient Joint View},
%             immediately AFTER the discussion paragraph of
%             Corollary~\ref{cor:compound-effect-binding} ("The corollary
%             follows by using $B$ as the separating authority ...") and
%             BEFORE Proposition~\ref{prop:ambiguous-cell}.
%   Rationale: T3.1 depends on thm:representation-insufficiency and on
%             cor:compound-effect-binding (it is their two-principle
%             strengthening), and its permissiveness-side failure mode feeds
%             directly into the ambiguous-cell lower bound that follows.
%             This keeps both propositions in the opening part of the
%             Security Analysis, as intended ("Security Analysis 开头").
%   Part 2 — related-work contrast paragraph:
%   File    : sections/related_work.tex
%   Position: at the end of \paragraph{Authorization, effects, and contracts},
%             immediately after the sentence "We build on complete mediation,
%             capabilities, confused-deputy analysis, information flow, and
%             provenance ~\cite{...}." (currently the last sentence of that
%             paragraph block).
%
% ----------------------------------------------------------------------------
% Dependencies (all existing; no new external citations needed)
%   used   : def:authorization-equivalence, thm:representation-insufficiency,
%            cor:compound-effect-binding, prop:authorization-quotient,
%            prop:ambiguous-cell; obligation O3 (forward-referenced: O1--O5
%            are formalized later in the same section, in "Conditional Effect
%            Confinement" — the draft uses "O3 (formalized below)" once).
%   defined: prop:principle-separation-effects, prop:whole-call-sufficient.
%   Cross-draft dependency: the counterexample paragraph cites
%            Table~\ref{tab:policy-family-sensitivity}, which resolves only
%            after e2_policy_family_table.tex is landed (frozen E2 numbers).
%
% ----------------------------------------------------------------------------
% Number policy (数字策略)
%   Only frozen-domain numbers appear (E2 power-set deletion counts on the
%   56-call domain: 16/16/16/8/4; common-contract baseline 118; typed 0).
%   No v17 and no V0-V3 numbers are referenced.
% ============================================================================

% ============================ PART 1: propositions ===========================

A natural question is whether classical enforcement principles already cover
this obligation. Two of them constrain the scope and timing of decisions but
not the representation on which decisions operate: least privilege restricts
how much authority a bound may instantiate, and complete mediation---the
obligation O3 formalized below---requires a guard decision before every
externally effectful transition. Neither constrains what the guard looks at.

\begin{proposition}[Least privilege and complete mediation do not imply safety]
\label{prop:principle-separation-effects}
Let $D$ be a finite instance domain and let $\rho_E$ be a monitor's
effect-side representation. Suppose the monitor satisfies least privilege
(bounds instantiate only the minimum necessary authority) and complete
mediation (every externally effectful transition has a preceding guard
decision, obligation O3). If $\rho_E$ is not authorization-sufficient on
$D$---in particular, if there exist $u,v\in D$ with $\rho_E(u)=\rho_E(v)$ and
some $B\in\mathcal{B}$ with $I_B(u)\neq I_B(v)$---then the monitor cannot be
both sound and permissive on $D$: it must either allow the unauthorized
context or withhold the authorized one.
\end{proposition}

\begin{proof}[Proof skeleton]
\begin{enumerate}
  \item By insufficiency, take $u,v\in D$ with $\rho_E(u)=\rho_E(v)$ and,
        without loss of generality, $I_B(u)=1$ and $I_B(v)=0$ for some
        $B\in\mathcal{B}$.
  \item Under the same complete view of $B$, the two authorization instances
        induce equal joint views, $R(u)=R(v)$, while $I(u)\neq I(v)$.
  \item Least privilege and complete mediation constrain the scope and timing
        of the decision, not the decision surface: any deterministic monitor
        $M$ satisfying both remains a function of $R$. By
        Theorem~\ref{thm:representation-insufficiency}, $M$ returns the same
        decision on $R(u)=R(v)$.
  \item Returning $\Allow$ violates soundness on $v$; returning $\Deny$ or
        $\Abstain$ violates permissiveness on $u$, incurring the
        Proposition~\ref{prop:ambiguous-cell} lower bound. The dichotomy is
        unavoidable. \qedhere
\end{enumerate}
\end{proof}

The proposition is the two-principle strengthening of
Corollary~\ref{cor:compound-effect-binding}: that corollary exhibits the
dichotomy for merged compound effects; the proposition shows that the
dichotomy persists even for monitors that satisfy least privilege and complete
mediation. The two principles are necessary obligations behind
Theorem~\ref{thm:single-commit} (via O3) but are not a substitute for
representation sufficiency.

\paragraph{Counterexample family.}
The separating structure is the calendar case of
Section~\ref{sec:security-analysis}: two calls execute the same effect
signature---\texttt{add\_participant(calendar event 6, target=}$\cdot$\texttt{)}---while
an authorization-relevant attribute differs; the monitor's representation
drops the attribute. Taking $u$ with target \texttt{internal@example.com}
($I_B(u)=1$ for a bound admitting internal members) and $v$ with target
\texttt{external@example.com} ($I_B(v)=0$), a monitor whose $\rho_E$ drops
the subject attribute satisfies both principles yet cannot separate $u$ from
$v$: allowing admits the unauthorized external invitee, while withholding
false-denies the authorized internal one. This family is instantiated by the
frozen finite domain: under the power-set authority family, deleting the
subject, date, or recurrence qualifier creates 16 authorization-separating
pairs each (payload 8, visibility 4), the common-field representation retains
118 such pairs, and the reviewed typed representation retains none
(Table~\ref{tab:policy-family-sensitivity}). A two-principle monitor over the
coarser representation inherits exactly these collisions.

\begin{proposition}[Existence of authorization-sufficient whole-call representations]
\label{prop:whole-call-sufficient}
For every finite domain $D\subseteq\mathcal{U}$, the map
\[
  \rho_W(u)=[u]_{\equiv_{\mathcal{B}}},
\]
assigning each context a canonical identifier of its
authorization-equivalence class, is an authorization-sufficient
representation on $D$, and it is a whole-call representation: one token per
context, not a sequence of effect-instance atoms. Consequently, least
privilege and complete mediation do not imply that an atom-level
representation is required for sound enforcement.
\end{proposition}

\begin{proof}[Proof skeleton]
\begin{enumerate}
  \item $\rho_W(u)=\rho_W(v)$ iff $u\equiv_{\mathcal{B}}v$, hence
        $I_B(u)=I_B(v)$ for every $B\in\mathcal{B}$: sufficiency holds by
        Definition~\ref{def:authorization-equivalence}.
  \item $\rho_W$ emits one class identifier per context, so it is a
        whole-call representation.
  \item By Proposition~\ref{prop:authorization-quotient}, the quotient is the
        coarsest sufficient representation; the cells of $\rho_W$ are exactly
        the quotient classes, so $\rho_W$ is its canonical implementation.
  \item Combined with Proposition~\ref{prop:principle-separation-effects}:
        the two principles neither imply safety (when the representation is
        insufficient) nor imply atom-level binding ($\rho_W$ is a sufficient
        whole-call counterexample). Atom-level binding is an engineering
        choice within the family of sufficient representations---supported by
        finite verification, registered-refinement convergence, and
        multi-family sensitivity evidence---not a logical consequence of the
        two principles. \qedhere
\end{enumerate}
\end{proof}

The existence claim carries no computability commitment: $\rho_W$ requires an
oracle for authorization equivalence, and the engineering objects that
approximate it are collision-complete finite validation and registered
refinement, not the map itself. The proposition's subject is the logical
implication relation between principles and representation granularity, not
the existence of a practical whole-call monitor.

% ======================= PART 2: related-work paragraph ======================
% Placement: sections/related_work.tex, end of the \paragraph{Authorization,
% effects, and contracts} block (see header note). Draft text:

% Least privilege and complete mediation constrain the scope and timing of
% enforcement, not the representation on which it operates.
% Proposition~\ref{prop:principle-separation-effects} constructs monitors that
% satisfy both principles yet remain unsafe whenever the effect-side
% representation merges authorization-inequivalent calls;
% Proposition~\ref{prop:whole-call-sufficient} exhibits an
% authorization-sufficient whole-call representation, so the two principles do
% not imply atom-level binding. Tool-effect binding is therefore not a
% relabeling of least privilege or of argument provenance: it is a property of
% the representation itself.

% ----------------------------------------------------------------------------
% NON-COMMITMENT LIST (must survive porting, T3 md §4):
%   no computable implementation of rho_W is claimed; no safety statement for
%   obligation combinations beyond the two principles; no open-domain
%   representation sufficiency.
% ============================================================================
